\section{Exact finite spectra for arbitrary hyperbolic matrices}
\label{sec:general-spectra}
The following finite character argument is independent of the Gaussian and
normal-form arguments below. In particular it does not assume a split torus
or a prime Hilbert-space dimension. This section reconciles the positive- and
negative-trace notes using one sign convention.

\subsection{A fixed-point-free cyclic action}
Let $N\geq3$ and $q\geq3$ be odd integers, and suppose
$A\in\mathrm{SL}(2,\mathbb Z/N\mathbb Z)$ satisfies
\begin{equation}\label{eq:general-free}
 A^q=I,\qquad \det(A^j-I)\in(\mathbb Z/N\mathbb Z)^\times
                      \quad(0<j<q).
\end{equation}
Use modular-centred Weyl operators $W(v)$ and $R\psi(x)=\psi(-x)$.
The assignment $W(v)\mapsto W(Av)$ is a $*$-automorphism of the full matrix
algebra, since it preserves the Weyl product and adjoint. It is implemented
by a unitary $U$, unique up to a scalar; this follows by mapping a complete
system of matrix units. No field hypothesis is required.

\begin{theorem}[Uniform parity sector]\label{thm:uniform-sector}
Under \eqref{eq:general-free}, $N\equiv c\pmod q$ for exactly one
$c\in\{1,-1\}$. A scalar normalization of $U$ gives $U^q=I$ and
\[
 \operatorname{Tr}(RU^j)=1\quad\hbox{for every }j,\qquad
 \operatorname{Tr}U^j=c\quad(q\nmid j).
\]
On $\mathcal H_c=\operatorname{ran}((I-cR)/2)$, every $q$th root has the same
multiplicity $n=(N-c)/(2q)$, and $\dim\mathcal H_c=(N-c)/2$.
On the complementary parity sector the multiplicities are $n$ at nontrivial
roots and $n+c$ at the root $1$.
\end{theorem}
\begin{proof}
Conjugation by $U^j$ permutes the Weyl basis. Its trace is consequently the
number of fixed Weyl indices, while the trace of conjugation by a unitary V
is $|\operatorname{Tr}V|^2$. Thus
\[
 |\operatorname{Tr}U^j|^2=1\quad(0<j<q),\qquad
 |\operatorname{Tr}(RU^j)|^2=1\quad\hbox{for all }j.
\]
For the second identity, a vector with $A^jv=-v$ is zero: applying the
relation $q$ times gives $v=-v$ and $2$ is a unit. The same argument modulo
each prime divisor of N proves invertibility, including prime powers.
The operators $RUR$ and U implement the same covariance. Their scalar ratio
is one because $\operatorname{Tr}U\ne0$. Normalize $U^q=I$.

Let $b_l$ be the difference between the even and odd multiplicities of the
$l$th root. Fourier inversion for the virtual cyclic character
$\operatorname{Tr}(RU^j)$ gives $b_l\in\mathbb Z$,
$\sum_l b_l^2=1$, and $\sum_l b_l=\operatorname{Tr}R=1$.
Exactly one coefficient is one. A further root-of-unity scalar puts it at
$l=0$. Hence $\operatorname{Tr}(RU^j)=1$ for every j.
For the full multiplicities $m_l$, Parseval yields
\[
 \sum_lm_l=N,\qquad \sum_lm_l^2=(N^2+q-1)/q.
\]
Write $N=qk+r$, $0\leq r<q$. The least sum of squares of q integers with
this sum is $qk^2+2kr+r$: moving a pair which differs by two or more closer
together reduces its squared sum. It follows that $r(q-r)\leq q-1$.
Integrality also forces $r^2\equiv1\pmod q$, so $r\ne0$; therefore
$r=1$ or $q-1$. For either residue the squared-sum minimum is attained.
All $m_l$ equal $(N-c)/q$ except one, which differs by c. Since $(N-c)/q$
is even and $b_0=1$, that exception must be at zero. Splitting by the already
known parity differences gives the claimed formulas.
\end{proof}

\subsection{An explicit sequence, including negative trace}
Let now $A\in\mathrm{SL}(2,\mathbb Z)$ be hyperbolic, put
$a=|\operatorname{tr}A|\geq3$, $s=\operatorname{sgn}(\operatorname{tr}A)$,
and let $u_0=0,u_1=1,u_{j+1}=a u_j-u_{j-1}$. For every odd prime $K>3$
not dividing $a-2s$, set
\begin{equation}\label{eq:general-moduli}
 m=(K-1)/2,\qquad N_K=u_{m+1}+s u_m,\qquad
 c_K=\left(\frac{a-2s}{K}\right).
\end{equation}
Here the last expression is a Legendre symbol, not an eigenvalue of A.
Then $N_K$ is odd and
\begin{equation}\label{eq:general-identity}
 A^K-I=s^mN_K A^m(A-I).
\end{equation}
Moreover $\gcd(N_K,a-2s)=1$, so Theorem~\ref{thm:uniform-sector} applies
with $q=K$ and sign $c=c_K$.

Here are details of the signed extension. Put $B=sA$, of trace a. Summing
$A^j=s^jB^j$ symmetrically around m and using Cayley--Hamilton yields
$\sum_{j=0}^{K-1}A^j=s^mN_KA^m$, proving
\eqref{eq:general-identity}. For $s=1$, the polynomial $N_K(a)$ evaluates
to K at a=2. For $s=-1$, it evaluates to $(-1)^mK$ at a=-2.
These evaluations prove the gcd assertion under the stated exclusion.
Modulo 2, the recurrence is alternating when a is even and has period three
when a is odd; the only obstruction is $K=3$, already excluded.
If $A^jv=v$ modulo a prime dividing $N_K$ with $0<j<K$, Bezout applied
to j,K and \eqref{eq:general-identity} gives $Av=v$. Invertibility of $A-I$
then gives $v=0$. Thus all required determinants are units.

The sign can be evaluated without a square root modulo $N_K$. For positive
trace the integer polynomial identity
\[
 N_K(t)=2^{1-K}\sum_{\substack{1\leq j\leq K\\j\text{ odd}}}
 \binom Kj(t+2)^{(K-j)/2}(t-2)^{(j-1)/2}
\]
gives $N_K(t)\equiv(t-2)^{(K-1)/2}\pmod K$. For negative trace apply
it at $t=-a$ and use $N_K^{(+)}(-a)=(-1)^mN_K^{(-)}(a)$.
Euler's criterion proves \eqref{eq:general-moduli}. With
$\lambda=(a+\sqrt{a^2-4})/2$ one also obtains
\[
 N_K=\frac{\lambda^{K/2}-s\lambda^{-K/2}}{\sqrt{a-2s}},
 \qquad K\asymp\log N_K.
\]
For trace four and $K=11$, $N_K=989=23\cdot43$, $c=-1$:
the uniform sector is even, of dimension 495 and multiplicity 45.
The discriminant 12 is a square modulo 23 and a nonsquare modulo 43.
For trace $-6$ and $K=7$, $N_K=169=13^2$, $c=1$:
the uniform odd sector has dimension 84 and multiplicity 12. These examples
require neither a split characteristic polynomial nor squarefreeness.

\subsection{The characteristic for general integer covariance}
The physical lift $Q_N(a,b)=(-1)^{ab}W_N(a,b)$ used for T is not invariant
under every integer symplectic matrix. For general A use instead
\[
 Q_N^*(a,b)=(-1)^{ab+a+b}W_N(a,b).
\]
The function $ab+a+b$ modulo two equals one on all three nonzero vectors,
so it is invariant under $\mathrm{SL}(2,\mathbb F_2)$. Therefore
\[
 Q_N^*(v)Q_N^*(w)=e^{\pi i\omega(v,w)/N}Q_N^*(v+w),\qquad
 U_NQ_N^*(v)U_N^*=Q_N^*(Av).
\]
This is an explicit quantization convention, not a claim that the old and
new finite matrices coincide. The additional sign represents a fixed
half-torus translation of symbols. Parity still corresponds to $x\mapsto-x$.
Whether the locality and return-word arguments later in the paper transfer
to this characteristic and a given $A$ depends on their remaining inputs,
which are not checked here. What is verified here is exactly the displayed
Weyl product law with the same cocycle, exact $A$-covariance for every
integer $A$, and the parity correspondence. The transfer is therefore
conditional on (i) finite support towers for the $A$-orbit with the gap and
noncancellation properties those arguments assume, and (ii) the finite-order
and unit-determinant hypotheses on $A$ modulo $N$, in the form of
\eqref{eq:general-free}, wherever they are used. No verbatim transfer to an
arbitrary such $A$ is asserted. The associated classical map is
$S=A^{\mathsf T}$.
